\chapter{图像生成与编辑}
\label{chap:vision-image-generation}

“画一只玩具狗”和“把这只玩具狗旁边的杯子换成花瓶”都有许多合理结果，但需要保持的内容不同。前者主要服从创作条件，后者已有一幅原图，玩具狗的身份和未指定背景也属于约束。若要求根据绕玩具狗拍摄的照片生成背面视图，还要说明哪些部分受到实际观测支持。本章按这些变化关系，讨论新图像的创建、已有图像的转换、指定内容的编辑和观测场景的新视角合成。

生成模型的通用训练与采样见\BookChapterRef{chap:pgm-generative}{第二卷“模型”中的“生成模型”章}，图文表示和媒体条件生成接口将在\ref{chap:vision-multimodal-models}章集中说明。这里保留使用这些接口所必需的输入、条件、监督与输出过程，把重点放在图像任务特有的要求上。贯穿示例中的候选、分数和相机数值均为教学构造，不表示模型实测结果。

\section{新图像的创建}
\label{sec:vision-generation-creation}

创建新图像时，需要先确定要求，再选择能够表达这些要求的条件，最后生成并筛选候选。这些是创作过程的职责，不要求模型内部恰好存在三个同名模块。自由采样可以只规定分辨率和大致图像分布，有条件创建则还要核验条件是否真正进入结果。

\subsection{内容与成果要求的确定}
\label{sec:vision-generation-requirements}

考虑“红色玩具狗放在蓝色杯子左边”。其中至少有两个对象、各自的颜色、数量及左右关系。把它写成可检查记录，可以得到“一个玩具狗，红色；一个杯子，蓝色；狗在杯子左边”。这份记录供结果核验，不意味着生成器内部已经显式保存同样的对象表。

对象出现与属性绑定是两层要求。画面同时存在红、蓝两色，不能说明颜色分别属于正确对象；包含狗与杯子，也不能说明数量和空间关系满足。若还要求杯子上写“OPEN”，字符应作为单独条件检查，不能仅看整体图文相似度。

输出用途也会改变要求。插图允许多种合理材质和光照，产品展示可能需要固定主体形状，排版素材还要规定留白与裁剪范围。条件没有指定的部分可以变化，已经明确的部分则不能由“画面自然”替代。对多个可接受结果，应保留变化空间，而非把某一参考图的全部像素当作唯一答案。

\subsection{创作条件的组织}
\label{sec:vision-generation-conditions}

文字容易描述语义，空间图更直接地规定位置，主体参考则提供某个具体玩具的外观。它们约束不同对象，也可能相互冲突。选择接口时，应说明每项条件如何编码、在哪一步影响生成，以及冲突时如何判断结果。

\subsubsection{潜空间条件采样与unCLIP的应用入口}
\label{sec:vision-generation-latent-entry}

潜空间生成先在压缩后的图像表示中采样，再由图像解码器还原像素。使用文字条件时，文字编码器产生词元特征，去噪网络在采样过程中读取这些特征；随机初值、采样器、步数和引导强度共同决定一次输出。图像解码器负责把最终潜变量变成$\bm{X}\in\mathbb{R}^{h\times w\times3}$，它不会自动验证文字中的数量和位置\scite{pgm-generative-rombach2022}

unCLIP将CLIP图像表示作为图像生成的条件。由文字创建时，先验模型根据文字产生可能的CLIP图像向量，再由生成模型形成图像；给定参考图像时，也可以直接编码参考，绕过文字到图像表示的先验，生成具有相关语义的变体。两条路径共享的是图像表示接口，不是必须经过同一个输入过程\scite{vision-unclip2022}

这种全局向量压缩了很多信息。参考中的玩具狗与新结果可能同属相近外观，但精确耳朵形状、徽标和原背景不一定被完整保留。若需要逐区域保持，应增加空间或局部条件；若要求同一主体的新姿态，还需要针对身份进行适配和评价。

实际使用时固定模型、提示、分辨率与总候选数，例如对同一描述生成四个随机种子的结果，比较姿态、背景和细节。相同种子只有在模型、采样过程及实现条件一致时才具有可重复含义。提高步数、改提示或增加候选都会改变总预算，不能只呈现最终选中图片而省略这些条件。

\subsubsection{ControlNet的空间条件支路}
\label{sec:vision-generation-controlnet}

文字“抬起左前爪”不容易精确指定每个关键点的位置。ControlNet把边缘、深度、分割或关键点图作为空间条件接入已有生成模型，使输入结构在去噪中提供约束\scite{vision-controlnet2023}
预训练主网络保留冻结参数，可训练支路由相应编码及中间模块复制而来，连接处使用零初始化卷积。

用一个模块说明其计算。设预训练模块为$F_{\bm{\theta}}$，复制支路为$F_{\bm{\theta}_{\mathrm{c}}}$，输入特征为$\bm{h}$，条件编码为$\bm{c}$，零初始化卷积为$Z_1,Z_2$，则示意输出为
\begin{equation}
 \bm{h}'=F_{\bm{\theta}}(\bm{h})+
 Z_2\!\left(F_{\bm{\theta}_{\mathrm{c}}}
       \bigl(\bm{h}+Z_1(\bm{c})\bigr)\right)\text{。}
 \label{eq:vision-controlnet}
\end{equation}
训练开始时，$Z_2$的输出为零，支路不扰动原模块结果。零输出不意味着整个支路永远没有学习信号：外层卷积的参数可先从其输入特征得到梯度，变为非零后再向复制模块传递梯度。图\ref{fig:vision-controlnet}用加法节点标明两条路径的关系。

\begin{figure}[htbp]
 \centering
 % Single-module sketch of ControlNet Eq.(2); no complete U-Net is implied.
\begin{tikzpicture}[x=1mm,y=1mm,font=\figurefont\small,
 b/.style={draw=BookBlue,fill=BookBlue!10,line width=.8pt,
  minimum width=16mm,minimum height=9mm,inner sep=1mm,align=center},
 op/.style={circle,draw=BookInk,fill=white,line width=.8pt,
  minimum size=5mm,inner sep=0},
 a/.style={-{Stealth[length=2mm]},BookInk,line width=.9pt,
  rounded corners=1mm}]
 \node at (5,0) (h) {$\bm{h}$};
 \node[b,draw=BookMuted,fill=BookPaper,minimum width=25mm] at (42,0)
  (base) {$F_{\bm{\theta}}$};
 \node[op] at (82,0) (sum) {$+$};
 \node at (101,0) (out) {$\bm{h}'$};
 \node[op] at (30,-25) (cond) {$+$};
 \node[b,minimum width=23mm] at (55,-25) (copy) {$F_{\bm{\theta}_{\mathrm{c}}}$};
 \node[b,minimum width=12mm] at (82,-25) (z2) {$Z_2$};
 \node at (5,-45) (c) {$\bm{c}$};
 \node[b,minimum width=12mm] at (30,-45) (z1) {$Z_1$};
 \draw[a] (h)--(base);
 \draw[a] (base)--(sum);
 \draw[a] (sum)--(out);
 \draw[BookInk,line width=.9pt] (14,0)--(14,-25);
 \draw[a] (14,-25)--(cond);
 \fill[BookInk] (14,0) circle (.6mm);
 \draw[a] (cond)--(copy);
 \draw[a] (copy)--(z2);
 \draw[a] (z2)--(sum);
 \draw[a] (c)--(z1);
 \draw[a] (z1)--(cond);
 \node[anchor=south,text=BookMuted] at (42,7) {冻结};
 \node[anchor=north,text=BookBlue] at (55,-32) {可训练副本};
 \node[anchor=north,text=BookBlue] at (82,-32) {从零连接};
\end{tikzpicture}

 \caption{ControlNet的模块级教学示意。主模块参数$\bm{\theta}$冻结，复制模块及零初始化连接可训练；相同输入进入两侧，条件只通过支路注入。零初始化让控制输出从零开始，训练后经加法调节主路径。完整网络在多个尺度接入控制特征，图中仅展开一处。}
 \label{fig:vision-controlnet}
\end{figure}

训练样本包含目标图像、文字和空间条件。可以从目标图像提取边缘或深度形成配对，给目标的潜变量加噪，再以条件噪声预测损失更新支路与连接。空间条件需要经过编码并对齐相应特征分辨率，不能直接把全尺寸深度图加到任意通道张量上。不同控制类型通常具有各自训练设置，使用时应加载与条件类型匹配的权重。

推理同时给出文字和控制图，在每次去噪中注入相应支路特征；控制强度可对这些修正作缩放。固定玩具狗姿态图，改变文字可探索材质与背景，改变强度则改变对输入结构的依从程度。控制强度为零使这条修正消失，强度增加也不保证任意条件都能满足。

多条控制支路可以把修正相加，但若深度图规定前爪在后、姿态图又规定前爪伸向相机，两种条件并没有一致的几何解。应核验条件提取器本身，并记录模型版本、条件分辨率和选择预算。边缘相近不等于主体身份相同，空间控制还需要与下一节的主体要求配合。

\subsubsection{DreamBooth的主体绑定与类别先验保持}
\label{sec:vision-generation-dreambooth}

普通类别词“玩具狗”无法区分某个具体玩具。DreamBooth用少量同一主体的参考照片和带特殊标识的类别提示微调生成模型，让标识与该主体外观建立联系。例如提示可包含一个选定的稀有标识和“toy dog”，使用时再接“在雪地中”。这个标识是文字条件的一部分，不是逐像素复制指令\scite{vision-dreambooth2023}

只在少量主体照片上训练，模型可能连参考背景一起记住，还可能把一般的“玩具狗”类别词也缩窄成这一只。原方法由微调前模型生成一般类别样本，并用类别先验保持目标约束这一变化。沿用噪声预测记法，可将训练目标概括为
\begin{equation}
 L=L_{\mathrm{subj}}+\lambda_{\mathrm{pr}}L_{\mathrm{class}}\text{，}
 \label{eq:vision-dreambooth}
\end{equation}
其中$L_{\mathrm{subj}}$在主体图像及“特殊标识+类别名”条件下预测所加噪声，$L_{\mathrm{class}}$在原模型生成的一般类别图像与普通类别提示下执行相应去噪训练。若模型直接预测清晰图像而非噪声，需使用与该参数化匹配的损失权重，不能不加转换地混用两种目标。

训练后保留微调参数，输入“该标识的玩具狗在雪地中”便可采样，不再需要同时提供目标雪地图像。类别样本约束普通类别语义，主体样本约束具体外观，两者提供不同训练信息。增大类别项并不会补齐主体照片从未展示的背面，也不能保证特殊标识在所有姿态下都保持细节。

核验时应把新结果与多张参考并排比较耳朵、斑纹和形状，另看是否只是复制了原背景。如果只有正面参考，生成背面的合理纹理仍可能来自类别先验。多个主体各自保持又增加了标识绑定与相互遮挡问题，不能从单主体成功直接推出多主体能力。

\subsection{图像生成与候选选择}
\label{sec:vision-generation-selection}

给定相同条件，采样过程可以主动改变生成方向，生成之后也可以从候选中挑选。这两种操作都可能使用图文相似度，却在不同时间、对不同对象起作用，成本也不同。

\Needspace{60mm}
\subsubsection{GLIDE的CLIP过程引导与无分类器引导对照}
\label{sec:vision-generation-glide-guidance}

GLIDE比较了噪声感知CLIP引导与无分类器引导。前者训练可读取当前带噪图像及噪声步的图像编码器，让文字相似度的梯度在每步修正采样；它使用的不是只在清晰图片上训练的原始CLIP\scite{vision-glide2022}

设当前反向转移的均值与协方差为$\bm{\mu}^{(t)},\bm{\Sigma}^{(t)}$，噪声感知图文相似度为$s(\bm{X}^{(t)},c,t)$，则均值修正具有形式
\begin{equation}
 \tilde{\bm{\mu}}^{(t)}
 =\bm{\mu}^{(t)}
  +a\,\bm{\Sigma}^{(t)}
   \nabla_{\bm{X}^{(t)}}s(\bm{X}^{(t)},c,t)\text{。}
 \label{eq:vision-clip-process-guidance}
\end{equation}
$a$为引导强度，梯度针对当前带噪图像，目的在于把后续采样推向文字条件。过程需要额外图文网络与梯度计算；将普通CLIP直接接入相同位置，是另一种分布和方法设置，不能沿用GLIDE对噪声感知模型的比较结论。

无分类器引导不运行独立分类器。训练以一定概率去掉文字，使同一去噪网络能分别预测有条件与无条件噪声；使用时按
$\hat{\bm{\epsilon}}=\bm{\epsilon}_{\varnothing}+b(\bm{\epsilon}_{c}-\bm{\epsilon}_{\varnothing})$
组合，$b=1$退回普通条件预测，$b>1$沿条件差异外推。它仍需要约定的条件与无条件预测计算，并不是没有额外推理成本。

固定玩具狗提示与噪声种子改变引导强度，可以观察主体、背景和多样性怎样变化。更强引导可能使某些条件更突出，也可能改变色彩和构图。对“两个杯子”“蓝色杯子在右侧”等组合条件，仍要逐项核验，不能把相似度上升解释为每项要求都更正确。

\subsubsection{图文相似度重排序与条件筛选}
\label{sec:vision-generation-reranking}

重排序发生在候选图像已经形成之后。对$k$幅图像$\bm{X}_1,\ldots,\bm{X}_k$，计算单位图文向量的相似度$s_i$并排序，它不会修改候选像素，也不能创造候选中不存在的正确解。GLIDE对照中的DALL-E使用了这类生成后CLIP筛选；不同对照还使用了不同总候选量，应连同选择协议阅读。

表\ref{tab:vision-generation-candidates}给出四个构造候选。若仅取最高全局相似度，会选择数量错误的A；若先筛选数量、属性归属与关系，再按约定标准排序，则只有D满足全部要求。这里的分项检查来自人工或独立检测、读取工具，不是CLIP分数已经提供了这些保证。

\begin{table}[htbp]
 \centering
 \small
 \begin{tabularx}{\linewidth}{@{}c c X c@{}}
  \toprule
  候选 & 构造相似度 & 按请求核验的主要结果 & 保留\\
  \midrule
  A & 0.34 & 出现两只玩具狗 & 否\\
  B & 0.32 & 狗为蓝色，杯子为红色 & 否\\
  C & 0.30 & 属性正确，但狗在右侧 & 否\\
  D & 0.28 & 数量、颜色归属与左右关系均满足 & 是\\
  \bottomrule
 \end{tabularx}
 \caption{全局相似度与逐项条件检查的教学对照。分数人为设置，用于展示两种选择规则的区别；不代表任何模型的实测排序。}
 \label{tab:vision-generation-candidates}
\end{table}

如果一个请求生成了100幅图，人工筛出1幅，报告时应保留100幅的生成成本和筛选规则。按最佳候选计分与按一次随机生成计分衡量的是不同系统。没有候选满足硬条件时，应返回条件未满足或继续按明确预算修改，而不是把最高分候选自动称为成功。

\section{已有图像的转换}
\label{sec:vision-generation-translation}

转换以已有图像为起点，获得另一种风格、呈现或成像形式。它的核心是建立输入与输出的对应关系；训练是否有配对参考，决定了哪些对应可直接监督。空间布局到照片也能使用同样的条件映射工具，本节借此解释机制，具体创作要求仍由\ref{sec:vision-generation-creation}节规定。

\subsection{转换关系与保持要求的规定}
\label{sec:vision-generation-translation-requirements}

把玩具狗照片转成绘画时，纹理和笔触可以改变，姿态和对象数量可能需要保持。把一种成像方式转换为另一种时，还可能要求每个位置与实际结构对应。两种用途即便采用相同网络，允许的误差也不同。

数据应与这种要求相符。对齐的照片—绘画配对可约束具体样本关系，独立的照片集与绘画集只说明两个域中分别有哪些外观。若把没有逐位置对应的图像直接作像素监督，损失会强迫网络拟合无意义的差异。

\subsection{转换映射的获得与使用}
\label{sec:vision-generation-translation-mapping}

\subsubsection{pix2pix的配对条件映射与局部判别}
\label{sec:vision-generation-pix2pix}

pix2pix使用对齐图像对$(\bm{X},\bm{Y})$学习条件映射。生成器$G$读取$\bm{X}$并输出目标形式，判别器$D$同时读取输入与候选输出，从而判断它们是否像真实配对。原方法的U形生成器用跳连把较细的位置特征传到解码端，局部判别器则在多个图像块上提供真假反馈\scite{vision-pix2pix2017}

以边缘图到玩具狗照片为例，$\bm{X}$提供轮廓，$\bm{Y}$提供该轮廓对应的颜色与纹理。标准条件对抗目标为
\begin{equation}
 \begin{aligned}
 L_{\mathrm{cGAN}}(G,D)
 &=\mathbb{E}_{\bm{X},\bm{Y}}[\log D(\bm{X},\bm{Y})]\\
 &\quad+\mathbb{E}_{\bm{X}}[
       \log(1-D(\bm{X},G(\bm{X})))],\\
 L_1(G)&=\mathbb{E}_{\bm{X},\bm{Y}}
       [\lVert\bm{Y}-G(\bm{X})\rVert_1]\text{。}
 \end{aligned}
 \label{eq:vision-pix2pix}
\end{equation}
训练交替更新判别器与生成器。判别器最大化上述对抗项；生成器在实际实现中改用最大化$\log D(\bm{X},G(\bm{X}))$的非饱和目标，并最小化加权像素误差$\lambda L_1$。这种替换在判别器容易辨认假图时提供更有用的梯度。这里省略随机变量，强调输入—目标对应；原工作发现显式噪声容易被忽略，最终在训练和使用时均保留dropout，但输出变化仍有限，不能据此把生成器当成具有充分多样性的条件采样器。

像素项把输出拉向这幅配对参考，局部判别项鼓励目标域的局部外观。两者都不会自行修正训练对中的错位。若每个真实轮廓总与向右偏移的照片配对，网络可能学到这种偏移；若多个合理纹理共享同一输入，像素目标也可能倾向平均。局部真实还不保证整幅图像的几何和身份正确。

使用时只调用生成器，不再需要判别器或目标参考。对于照片到另一成像形式的转换，应另外检查真实对应与诊断所需结构，不能把训练域中看似逼真的输出解释成新的测量观测。

\subsubsection{CycleGAN的无配对双向映射}
\label{sec:vision-generation-cyclegan}

只有照片域$A$与绘画域$B$的独立样本时，CycleGAN训练$G:A\to B$和$F:B\to A$两个方向，并为两个域分别设置判别器。目标域判别要求输出像该域图像，循环一致性则要求转换后还能回到输入\scite{vision-cyclegan2017}

循环损失为
\begin{equation}
 L_{\mathrm{cyc}}=
 \mathbb{E}_{\bm{x}\sim p_A}
 [\lVert F(G(\bm{x}))-\bm{x}\rVert_1]
 +
 \mathbb{E}_{\bm{y}\sim p_B}
 [\lVert G(F(\bm{y}))-\bm{y}\rVert_1]\text{。}
 \label{eq:vision-cycle-consistency}
\end{equation}
训练把两个域的对抗项和加权循环项结合。照片$\bm{x}$经$G$生成绘画，再经$F$返回照片，输入本身就是往返监督，无需知道它应对应哪幅真实绘画。需要额外颜色保持时，还可在特定设置中加入输入已属于目标域时少改动的约束，但它不等于逐样本跨域配对。

往返可恢复并不能唯一确定正确语义。假设两个域各只有一幅狗图和一幅猫图，且概率都为$1/2$。$G$把狗映成猫、猫映成狗，$F$执行逆映射，两个输出分布都正确，循环误差也为零，但对象语义完全交换。图\ref{fig:vision-cycle-ambiguity}将这个反例与保持类别的映射并排显示。

\begin{figure}[htbp]
 \centering
 % Both maps are bijections of two equiprobable categories.
% Target order is changed in the right panel to avoid line crossings.
\begin{tikzpicture}[x=1mm,y=1mm,font=\figurefont\small,
 b/.style={draw=BookBlue,fill=BookBlue!8,line width=.8pt,
  minimum width=12mm,minimum height=8mm,inner sep=1mm},
 a/.style={{Stealth[length=1.7mm]}-{Stealth[length=1.7mm]},
  BookInk,line width=.9pt}]
 \node[font=\figurefont\small\bfseries] at (20,14) {保持类别};
 \node[font=\figurefont\small\bfseries] at (79,14) {交换类别};
 \node[anchor=south,text=BookMuted] at (5,5) {域$A$};
 \node[anchor=south,text=BookMuted] at (35,5) {域$B$};
 \node[anchor=south,text=BookMuted] at (64,5) {域$A$};
 \node[anchor=south,text=BookMuted] at (94,5) {域$B$};
 \node[b] at (5,0) (ad) {狗};
 \node[b] at (35,0) (bd) {狗};
 \node[b] at (5,-17) (ac) {猫};
 \node[b] at (35,-17) (bc) {猫};
 \node[b] at (64,0) (ad2) {狗};
 \node[b] at (94,0) (bc2) {猫};
 \node[b] at (64,-17) (ac2) {猫};
 \node[b] at (94,-17) (bd2) {狗};
 \foreach \u/\v in {ad/bd,ac/bc,ad2/bc2,ac2/bd2} {
  \draw[a] (\u)--node[above=1mm] {$G/F$}(\v);
 }
 \node[anchor=north] at (20,-25) {循环误差为0};
 \node[anchor=north] at (79,-25) {循环误差也为0};
\end{tikzpicture}

 \caption{循环一致性不能唯一确定语义对应。两域各有等概率的狗、猫两项。左侧保持类别，右侧交换类别；每条往返路径都回到原项，且目标域边缘分布相同。反例是离散教学构造，用来说明循环与域判别约束仍有多解。}
 \label{fig:vision-cycle-ambiguity}
\end{figure}

使用照片到绘画方向时，只运行$G$；反向生成器和两个判别器是训练职责，不需要每次转换都执行完整循环。对玩具狗仍应检查姿态、身份和文字，对测量或医学成像则需要更强的配对及事实依据。无配对域转换可以形成视觉风格，不能凭循环重建为未观测结构提供真实性。

\Needspace{45mm}
\subsubsection{感知损失监督的固定风格转换}
\label{sec:vision-generation-perceptual-style}

若目标是一幅指定风格图的纹理，可以用固定视觉网络比较内容与风格，再训练一个前向转换网络\scite{vision-perceptual2016}
Johnson等人的方法区分图像变换网络与损失网络：前者生成图像并更新参数，后者提供特征度量并保持固定。

设变换网络输出$\widehat{\bm{X}}=G_{\bm{\theta}}(\bm{X})$，固定损失网络第$\ell$层的特征为$\bm{F}_\ell(\bm{X})\in\mathbb{R}^{c_\ell\times n_\ell}$。内容项比较对应位置特征，风格项比较通道之间的相关统计。令
\begin{equation}
 \bm{\Gamma}_\ell(\bm{X})
 =\frac{1}{c_\ell n_\ell}
   \bm{F}_\ell(\bm{X})\bm{F}_\ell(\bm{X})^\top\text{，}
 \label{eq:vision-style-gram}
\end{equation}
则可组合输出与输入的特征平方差，以及输出与固定风格图$\bm{S}$的
$\lVert\bm{\Gamma}_\ell(\widehat{\bm{X}})-\bm{\Gamma}_\ell(\bm{S})\rVert_{\mathrm{F}}^2$。归一化方式会影响损失权重，必须与采用的实现一致。

Gram矩阵对空间位置作求和，保留通道共同出现的统计，却不直接记录这些纹理出现在何处。例如单通道两位置特征$(1,0)$与$(0,1)$产生相同Gram值，位置已经交换。内容项因此承担保留较高层结构的职责，风格项则鼓励某种纹理统计；二者的权重决定输出怎样折中。

训练在许多内容图上更新$G_{\bm{\theta}}$，固定同一风格图，之后使用时只需一次生成器前向，不必对每张新图重新优化像素，也不再需要运行损失网络。这个原始固定风格路线不能直接解释成“推理时任意输入一幅风格图即可转换”。玩具狗的笔触可改变，徽标和细小字符是否保留仍需另查。

\subsection{转换结果与对应关系的核验}
\label{sec:vision-generation-translation-check}

转换结果应同时具有目标形式和规定的输入对应。照片仍像照片，说明风格可能未充分改变；绘画自然却把狗变成猫，则说明语义保持失败。把图像恢复成输入的循环误差很小，只能回答往返约束，不能代替中间结果的检查。

固定输入后，可沿对象、轮廓、文字与颜色逐项记录变化，按用途区分允许和不允许的改变。配对训练出现系统性错位时，优先检查数据；没有配对依据却要求严格成像保真时，需要增加相应监督或收窄用途。下一节的编辑将变化要求交给当前请求，而不是预先固定一个输入输出域关系。

\section{指定内容的编辑}
\label{sec:vision-generation-editing}

编辑给定原图，并通过当前请求规定变化。它可以改局部，也可以改全图，但都需要检查修改是否完成、应保留内容是否仍在。语言描述目标，掩码表达范围，参考图像提供期望外观，这几种输入的作用应分别说明。

\subsection{修改目标与范围的确定}
\label{sec:vision-generation-editing-scope}

“把玩具狗旁边的杯子换成花瓶”至少包含待替换对象、替换后的内容和周围保留物。若画面有两个杯子，仅给出对象类别还不足以确定操作范围。可以先用\ref{chap:vision-image-recognition}章的检测与分割取得候选，再依据关系或人工选择消歧；更复杂的表达定位将在\ref{chap:vision-multimodal-understanding}章展开。

沿用上一章的掩码约定，$\bm{M}=1$表示保留区，$1-\bm{M}$表示允许生成区。若花瓶比原杯子高，只遮住原杯子轮廓可能没有容纳新物体的空间，还可能留下杯子原来的阴影。合理范围应覆盖允许变化的接触区和阴影，并明确玩具狗是否可被遮挡。掩码是当前任务要求，不是生成器自动推断出的正确编辑边界。

\subsection{修改的实施与原内容保持}
\label{sec:vision-generation-editing-methods}

\subsubsection{GLIDE的图像与掩码条件修补}
\label{sec:vision-generation-glide-inpainting}

GLIDE的修补模型在当前带噪RGB之外，还读取未遮挡图像与掩码。训练时随机擦去图像区域，把剩余RGB和一个掩码通道作为额外条件；新增输入通道的权重从零开始，微调模型去噪恢复完整训练图像。于是模型在每步都能读取清晰的保留区域，而不只是当前带噪结果。

这一结构与“按普通生成模型采样，再覆盖已知像素”不同。后一做法可以在输出端强制保留某些像素，但没有因此让模型在训练中学会如何利用清晰上下文。GLIDE通过显式条件训练连接保留区域和目标内容；其上采样模型还读取完整低分辨率图像，以及高分辨率的未遮区域，输入职责也有所区别。

使用时给出原图、杯子及可变周边的掩码、花瓶描述与随机初值，反复去噪得到候选。应检查花瓶是否替换到正确位置，边缘、反光和投影是否与场景协调，玩具狗是否保持。若只要求恢复遮挡前的杯子，则属于\ref{sec:vision-restoration-inpainting}节的复原目标；这里的“花瓶”明确允许创造原图中不存在的内容。

掩码条件提供范围信息，但能否严格保持范围之外的每个像素还要看实现和后处理。硬性保留时可按已声明掩码回填原图，随后重新检查边界；这种回填是额外操作，不能把它造成的逐像素相等归为生成模型的能力。

\Needspace{60mm}
\subsubsection{InstructPix2Pix的合成编辑配对与双条件引导}
\label{sec:vision-generation-instruct-pix2pix}

用户有时只给一句指令，没有手工掩码。InstructPix2Pix学习从原图和编辑指令直接生成结果，监督数据采用合成三元组“原图、指令、编辑后图像”。语言模型根据原描述生成编辑指令与新描述，再用图像生成模型把两份描述转成对应图像\scite{vision-instruct-pix2pix2023}

独立生成两幅图会同时改变姿态、背景等大量内容，因而不适合作为局部编辑配对。原方法借助Prompt-to-Prompt思想，在两次生成中共享噪声和部分注意力过程，鼓励对应；其补充材料所述实现对所有编辑采用早期若干步的自注意力权重替换。随后按图文表示中的变化方向等规则筛选配对。共享计算降低了无关变化，却没有保证每个合成指令都被正确执行。

训练时将目标图像编码并加噪，得到$\bm{z}^{(t)}$，将原图编码为$\bm{z}_{I}$，两者沿通道拼接输入去噪网络，文字编码器则读取指令。噪声预测损失为
\begin{equation}
 L=\mathbb{E}\!\left[
 \left\lVert\bm{\epsilon}-
 \epsilon_{\bm{\theta}}
   \bigl(\bm{z}^{(t)},t,\bm{z}_I,c_T\bigr)
 \right\rVert_2^2\right]\text{。}
 \label{eq:vision-instruct-training}
\end{equation}
$c_T$是指令条件，$\bm{\epsilon}$是加到目标潜变量上的噪声。训练从预训练生成模型初始化，新增图像条件通道的权重初始化为零，并随机丢弃部分图像或文字条件，使模型学会不同条件组合下的预测。

使用时不需要目标图像，也不需要用户写出完整的新场景描述。输入原图与“把杯子换成花瓶”，从随机目标潜变量开始采样。记同一步的三种噪声预测为
$\bm{\epsilon}_0=\epsilon_{\bm{\theta}}(\bm{z}^{(t)},t,\varnothing,\varnothing)$、
$\bm{\epsilon}_I=\epsilon_{\bm{\theta}}(\bm{z}^{(t)},t,\bm{z}_I,\varnothing)$
和$\bm{\epsilon}_{I,T}=\epsilon_{\bm{\theta}}(\bm{z}^{(t)},t,\bm{z}_I,c_T)$，双条件引导为
\begin{equation}
 \hat{\bm{\epsilon}}
 =\bm{\epsilon}_0
  +s_I(\bm{\epsilon}_I-\bm{\epsilon}_0)
  +s_T(\bm{\epsilon}_{I,T}-\bm{\epsilon}_I)\text{。}
 \label{eq:vision-instruct-guidance}
\end{equation}
第二项强调图像条件，第三项在图像条件已给定时强调指令带来的变化。它不是把三个独立概率相加；这些向量是去噪计算的输出。

以两个分量示意，令$\bm{\epsilon}_0=(0.1,0.2)$、$\bm{\epsilon}_I=(0.2,0.1)$、$\bm{\epsilon}_{I,T}=(0.4,-0.1)$，取$s_I=1.5,s_T=5$，结果为$(1.25,-0.95)$。若两项强度均为1，代数上就恢复$\bm{\epsilon}_{I,T}$。改变强度会改变整条采样轨迹，不能从这个线性组合推出某个对象的修改量也线性变化。

较强图像条件通常有助于保留输入结构，较强指令条件会强化修改，但这种取舍需要对当前输入核验。原方法没有显式编辑掩码，因此未指定区域也可能改变。数量、空间关系和重复编辑的错误，还可能继承自合成数据及其筛选器。

\subsection{修改结果的核验与迭代}
\label{sec:vision-generation-editing-check}

结果比较应同时放入原图、编辑后图像、要求记录和允许变化区域。杯子仍在是目标未改，花瓶出现在另一侧是改错位置，玩具狗斑纹变化是保留失败，物体与桌面之间出现断边则属于边界问题。把这些错误混成一个“编辑质量”分数，不容易确定应修改哪一步。

对明确的保留掩码$\bm{M}$，可计算保留区平均变化
\begin{equation}
 E_{\mathrm{keep}}=
 \frac{\sum_{i,j,k}m_{i,j}
       |x'_{i,j,k}-x_{i,j,k}|}
 {c\sum_{i,j}m_{i,j}}\text{，}
 \label{eq:vision-edit-preservation}
\end{equation}
分母要求至少有一个保留像素。若整图照搬输入，$E_{\mathrm{keep}}=0$，但杯子没有变成花瓶；因此这项必须与修改完成度一起报告。允许全图色调变化时，逐像素保留误差也需要换成与要求相符的判断。

为了强制范围外保持，可合成$\bm{X}_{\mathrm{out}}=\bm{M}\odot\bm{X}+(1-\bm{M})\odot\bm{X}'$。它能保证所选保留像素相等，却可能把新花瓶的阴影切断。若允许羽化掩码，就要说明混合带也会改变，不能同时声称边界带严格不变。

迭代时应持续对照最初的原图和保留要求，避免每轮只与上一轮比较，逐步积累身份漂移。缩小范围、调整条件强度和重新生成都需记录轮数与候选预算；最终成功例应连同得到它的制作成本解释。

\section{观测场景的新视角合成}
\label{sec:vision-generation-novel-view}

\term{新视角合成}（novel view synthesis）根据已拍摄视图生成未拍摄视角的图像。这里的目标是呈现同一个观测场景，所需约束来自相机、跨视图对应和场景表示。若只用一幅图加生成先验构造可能的背面，观测支持更弱，相应空间内容创建将在\ref{chap:vision-multimodal-generation}章继续讨论。

\subsection{观测与相机条件的组织}
\label{sec:vision-generation-cameras}

绕静止玩具狗拍摄多张照片，记录每张图的内参、位置和朝向，并留出不参与场景拟合的测试视图。相机参数把二维像素连接到三维射线，几何估计可复用\ref{chap:vision-image-recognition}章的对应与重建方法。拍摄覆盖决定哪些表面有多视角证据；只有正面照片时，背面并没有因为模型能渲染而变成已观测区域。

原始NeRF和3D Gaussian Splatting都采用逐场景优化：为这一组图像拟合一个场景，再用任意给定相机渲染。换到另一只玩具通常需要重新建立表示。前馈场景预测则把跨场景规律放入预训练模型，使用成本和先验来源不同，不能混用“训练时间”的含义。

静态、曝光一致和相机准确是重要条件。玩具移动、光照变化或标定误差，会让同一空间点需要解释互相矛盾的像素。优化可能通过错误密度或外观来吸收矛盾，训练视图变好不一定意味着场景几何更准确。

\subsection{场景表示的拟合}
\label{sec:vision-generation-scene-fitting}

\Needspace{65mm}
\subsubsection{NeRF的射线采样与体渲染监督}
\label{sec:vision-generation-nerf}

\term{神经辐射场}（neural radiance field, NeRF）用网络表示空间中的体密度与随观察方向变化的颜色。位置$\bm{x}\in\mathbb{R}^3$决定非负密度$\sigma(\bm{x})$，位置和单位方向$\bm{d}$共同决定颜色$\bm{c}(\bm{x},\bm{d})\in[0,1]^3$。密度用于计算射线经过该处时有多少贡献会被截获，颜色则允许同一位置在不同方向呈现不同高光\scite{vision-nerf2020}

相机像素对应射线$\bm{r}(s)=\bm{o}+s\bm{d}$，其中$\bm{o}$是相机中心，$s$是沿射线的距离。给定近远界$s_{\mathrm{n}},s_{\mathrm{f}}$，体渲染颜色为
\begin{equation}
 \begin{aligned}
 \bm{C}(\bm{r})
 &=\int_{s_{\mathrm{n}}}^{s_{\mathrm{f}}}
    T(s)\sigma(\bm{r}(s))
    \bm{c}(\bm{r}(s),\bm{d})\,\mathrm{d}s,\\
 T(s)&=\exp\!\left(
     -\int_{s_{\mathrm{n}}}^{s}\sigma(\bm{r}(u))\,\mathrm{d}u
    \right)\text{。}
 \end{aligned}
 \label{eq:vision-nerf-continuous}
\end{equation}
$T(s)$是到达该位置之前的透射率，近处密度越大，远处颜色的可见贡献越小。若$s$有长度单位，$\sigma$具有其倒数单位，使指数无量纲；实际拟合常把场景归一化。有限远界之外的背景可另加$T(s_{\mathrm{f}})\bm{c}_{\mathrm{bg}}$。

数值计算沿射线选取$n$个有序样本，用每段近似常量的密度和颜色积分。设段长$\delta_i$，得到
\begin{equation}
 \begin{aligned}
 \alpha_i&=1-\exp(-\sigma_i\delta_i),&
 T_i&=\prod_{j<i}(1-\alpha_j),\\
 w_i&=T_i\alpha_i,&
 \widehat{\bm{C}}&=\sum_{i=1}^{n}w_i\bm{c}_i
       +T_{n+1}\bm{c}_{\mathrm{bg}}\text{。}
 \end{aligned}
 \label{eq:vision-nerf-discrete}
\end{equation}
空乘积取1。$w_i$包含“先前未被遮住”和“在当前段产生贡献”两项，故远处颜色不能只按自身$\alpha_i$相加。权重与剩余背景权重之和为1；若不加背景，射线穿过稀疏场景时颜色和可能不足。

取两段$\alpha_1=0.5,\alpha_2=0.75$，颜色分别为红$(1,0,0)$和蓝$(0,0,1)$，背景为白$(1,1,1)$。第一段权重0.5，第二段为$(1-0.5)0.75=0.375$，背景为$0.5\times0.25=0.125$，结果为$(0.625,0.125,0.5)$。如果忘记透射率，直接用0.5与0.75加权，就会重复计算被前景遮住的贡献。图\ref{fig:vision-ray-compositing}同时展示顺序和权重。

\begin{figure}[htbp]
 \centering
 % Two intervals, alpha .5 and .75; background transmittance .125.
\begin{tikzpicture}[x=1mm,y=1mm,font=\figurefont\small]
 \draw[BookMuted,line width=.9pt,->] (5,0)--(100,0) node[right] {$s$};
 \draw[BookInk,line width=.9pt] (5,-3)--(0,-6)--(0,6)--(5,3)--cycle;
 \node[anchor=south] at (4,9) {相机};
 \fill[BookBlue!12] (24,-7) rectangle (41,7);
 \draw[BookBlue,line width=.9pt] (24,-7) rectangle (41,7);
 \fill[BookTeal!15] (58,-7) rectangle (75,7);
 \draw[BookTeal,line width=.9pt] (58,-7) rectangle (75,7);
 \draw[BookMuted,line width=1pt] (91,-7)--(91,7);
 \node[anchor=south] at (32.5,8) {$\alpha_1=0.5$};
 \node[anchor=south] at (66.5,8) {$\alpha_2=0.75$};
 \node[anchor=north] at (32.5,-9) {$T_1=1$};
 \node[anchor=north] at (66.5,-9) {$T_2=0.5$};
 \node[anchor=south] at (91,8) {背景};
 \node[anchor=east] at (7,-31) {$w$};
 \fill[BookBlue!20] (10,-35) rectangle (50,-27);
 \fill[BookTeal!20] (50,-35) rectangle (80,-27);
 \fill[BookPaper] (80,-35) rectangle (90,-27);
 \draw[BookInk,line width=.7pt] (10,-35) rectangle (90,-27);
 \draw[BookInk,line width=.7pt] (50,-35)--(50,-27);
 \draw[BookInk,line width=.7pt] (80,-35)--(80,-27);
 \node at (30,-31) {$0.5$};
 \node at (65,-31) {$0.375$};
 \node[anchor=north] at (30,-37) {第一段};
 \node[anchor=north] at (65,-37) {第二段};
 \node[anchor=north] at (85,-37) {$0.125$};
\end{tikzpicture}

 \caption{两段射线的体渲染教学例。近处段先取得0.5的贡献权重，后段只能使用剩余0.5中的0.75，背景使用最终透射率0.125。下方长度条表示实际权重，三段和为1。颜色值人为设定，用于复算，图中形状不表示真实物体测量。}
 \label{fig:vision-ray-compositing}
\end{figure}

训练从观测图像抽取像素射线，经上述积分生成颜色，再与相应真实像素计算平方误差，对网络参数反向传播。原始NeRF用坐标的多频正余弦编码帮助表示细节，并分别限制密度只依赖位置、颜色还依赖方向。改变这些依赖会改变表示的多视图约束，不能把网络看成直接记忆每个训练像素编号。

空空间与已遮挡区域仍需网络查询，原方法因此使用粗细两级采样。粗网络在分段随机采样点上形成$w_i$，把权重归一化为沿射线的分段分布，再采更多点；细网络在粗点与新增点的并集上重新计算颜色。粗、细颜色都接受像素损失。粗权重全接近零时，需要数值稳定处理以形成有效采样分布；新增样本最终参与非均匀积分，不能只把高权重点的颜色作不带段长的平均。

渲染新相机时，重新为它的每个像素构造射线并查询网络，不是把旧图像中的某个标签搬过来。若一张图有$p$条射线，每条有$n$次网络查询，基本计算随$pn$增长，具体代价还取决于网络和采样层次。未被拍到的表面、镜面反射和错误相机都可能使新视图失败；图像损失本身没有唯一确定真实几何。

\subsubsection{3D Gaussian Splatting的投影、透明度合成与密度调整}
\label{sec:vision-generation-gaussians}

对每条射线反复查询网络的成本较高。3D Gaussian Splatting改用一组显式三维高斯表示场景，把它们投影到图像平面并合成像素。每个高斯保存位置$\bm{\mu}$、协方差$\bm{\Sigma}$、不透明度$o$和方向相关的颜色参数；原方法从相机估计得到的稀疏点云初始化这些对象，再逐场景优化\scite{vision-gaussian2023}

协方差控制高斯的方向与伸展范围。可写$\bm{\Sigma}=\bm{R}_{g}\bm{S}\bm{S}^{\top}\bm{R}_{g}^{\top}$，其中$\bm{R}_{g}$是旋转，$\bm{S}$含正尺度，以保持合法协方差。变换到目标相机后，在高斯中心附近线性化透视投影，二维协方差近似为
\begin{equation}
 \bm{\Sigma}_{2}
 =\bm{J}\bm{R}_{c}\bm{\Sigma}
    \bm{R}_{c}^{\top}\bm{J}^{\top}\text{，}
 \label{eq:vision-gaussian-projection}
\end{equation}
其中$\bm{R}_{c}$将世界方向转到相机坐标，$\bm{J}\in\mathbb{R}^{2\times3}$是该点的投影Jacobian。所得二维椭圆决定哪些像素受到这个高斯影响；这是局部投影近似，不能把三维高斯的透视投影无条件视为精确二维高斯。

若投影中心为$\bm{\mu}_{2}$，像素$\bm{u}$处的有效不透明度可概括为
\begin{equation}
 a_i(\bm{u})=o_i\exp\!\left[
 -\frac12(\bm{u}-\bm{\mu}_{2,i})^\top
 \bm{\Sigma}_{2,i}^{-1}(\bm{u}-\bm{\mu}_{2,i})
 \right]\text{。}
 \label{eq:vision-gaussian-alpha}
\end{equation}
实现还使用有限足迹、数值阈值及稳定处理。将覆盖该像素的高斯按近远顺序合成，颜色形式为
$\sum_i a_i\bm{c}_i\prod_{j<i}(1-a_j)$，必要时加剩余背景。它与式\eqref{eq:vision-nerf-discrete}具有相同的前向透明度合成结构，但$a_i$来自二维投影足迹及学习的不透明度，并非逐段密度积分的直接替换。

原方法按图像块组织高斯并排序，再在GPU上合成，降低空区域遍历与重复操作。排序使用高斯中心的深度来近似可见次序，没有对每个像素执行精确的表面求交。颜色采用球谐系数表示随方向变化的外观，即用一组定义在方向球面上的基函数组合颜色。拟合损失结合像素$L_1$与结构相似性差异，更新位置、尺度、旋转、不透明度和颜色参数。

固定数量的初始高斯未必能表达所有细轮廓。原方法依据投影位置的梯度等信号选择需要加密的区域，小尺度高斯可以克隆，覆盖过大的高斯可以分裂为多个较小对象，并删除贡献很弱或不合适的高斯。玩具耳缘若只有一个大高斯覆盖，局部颜色变化可能被混在一起；分裂提供更多自由度，之后仍需优化位置、颜色和不透明度。分裂动作本身不保证图像误差立刻降低，也不是测得了新的表面点。

新相机仍需重新投影、排序和合成。显式高斯减少了MLP查询，却增加了场景参数存储，像素被大量高斯覆盖时合成成本也会上升。速度比较要固定分辨率、场景规模和硬件。高斯集合没有天然的三角网格拓扑或物理材质，渲染自然不能直接推出资产适合碰撞、制造或精确编辑。

\subsection{目标视角的渲染与核验}
\label{sec:vision-generation-rendering-check}

保留视图应在场景拟合前确定，不用其像素误差调节该场景参数。给定保留相机，比较渲染与真实图像的PSNR、SSIM或LPIPS，同时检查对象轮廓、细节和反光。若相机标定使用了哪些图像信息，也应纳入协议；已知测试相机与未知相机的任务条件不同。

单个视图还可能遗漏不连续现象。沿连续相机路径渲染，可以观察遮挡变化、漂浮斑点和突然外观跳变。连续路径上没有真实图像的位置只能做一致性观察，不能计算缺失参考的真实误差。几何评价则需要点云、深度或其他测量依据，不能从连续且逼真的动画直接推出准确表面。

对于玩具狗背面，训练覆盖不足时，两种场景表示都可能把先验或错误拟合延伸过去。若用途只是某些可见视角的展示，应明确适用视角范围；若需要完整三维资产，则还需要背面采集与几何核验。动态场景的时间变化由\ref{chap:vision-multimodal-generation}章扩展。

\section{生成结果的评价与研究方向}
\label{sec:vision-generation-evaluation}

创建、转换、编辑和新视角合成有不同参照。自由创建主要看条件和分布中的变化，转换与编辑还要对照输入，新视角则要对照观测场景。比较模型时固定提示、参考、分辨率、候选数量和修订预算，才能知道差异来自模型还是选择过程。

\subsection{条件记录与对象关系核验}
\label{sec:vision-generation-condition-check}

回到玩具狗与杯子的记录，逐项检查存在、数量、颜色归属和左右关系，主体适配还要检查参考外观。可以用检测器统计对象，用区域特征比较主体，用文字读取器检查杯子标识，但这些检查器本身也可能漏检或误读。自动分数应说明输入和输出，必要时抽样人工复核，不把工具的错误都算到生成器头上。

整体图文相似度适合衡量粗粒度相关性，无法独立证明复杂条件全部成立。若只汇总一个平均分，应同时保留关键分项及失败率，例如对象数量正确率和主体保持情况。多个属性绑定错误即使只影响小区域，也可能使整项任务失败。

\subsection{修改范围与输入对应核验}
\label{sec:vision-generation-change-check}

编辑需要同时检查“该改的改了”与“该留的留了”。原样返回玩具狗和杯子可以取得最小图像距离，却没有执行指令；把杯子变成花瓶同时改变狗的斑纹，则完成了局部目标但破坏保留条件。两类错误应分别记录。

转换还应考虑目标形式本身允许的变化。照片变成绘画后，像素距离大是预期现象，应结合内容对应与风格要求解释LPIPS等距离。新视角中，同一对象因投影而移动也不是编辑漂移；评价必须比较对应目标相机，而非直接与原相机图像逐像素相减。

\subsection{观测支持、制作成本与研究问题}
\label{sec:vision-generation-costs}

一个可交付成果的成本包含生成、检查和修订。主体微调增加前期适配，空间控制增加条件制作，新视角增加采集、标定和场景拟合。一次渲染很快，不表示完成整个场景也只需一次前向；从百幅候选中选出一幅，也不表示单次生成稳定满足要求。

组合条件冲突、多主体绑定、文字和精确数量控制、连续编辑漂移，以及缺少视角时的观测忠实性，分别暴露不同约束的不足。研究比较应定位具体错误，而非将所有失败概括为“生成质量不够”。若把生成图像用于训练，还需在真实目标任务上验证收益，相关数据与迁移问题回到范式卷的知识利用讨论。

\section{本章小结}
\label{sec:vision-generation-summary}

同一只玩具狗在创建、转换、编辑和新视角合成中承担不同角色：它可能只是一个类别要求，也可能是必须保持的输入主体，或有多视角证据的真实场景。条件控制和主体适配提供不同约束，配对与循环监督建立不同对应，掩码和指令又分别规定可变范围与目标变化。评价需要跟随这些职责，不能仅凭自然外观或全局相似度判断任务完成。

新视角图像由场景表示、相机和渲染共同产生，遮挡与观测覆盖决定哪些内容受到支持。制作过程还要计入采样、筛选、适配与拟合成本。后续连续画面和多模态生成会增加时间、身份与同步约束，本章对条件、保持关系和观测证据的区分仍然适用。
